% Options for packages loaded elsewhere
\PassOptionsToPackage{unicode}{hyperref}
\PassOptionsToPackage{hyphens}{url}
\PassOptionsToPackage{dvipsnames,svgnames,x11names}{xcolor}
%
\documentclass[
]{article}
\usepackage{amsmath,amssymb}
\usepackage{iftex}
\ifPDFTeX
  \usepackage[T1]{fontenc}
  \usepackage[utf8]{inputenc}
  \usepackage{textcomp} % provide euro and other symbols
\else % if luatex or xetex
  \usepackage{unicode-math} % this also loads fontspec
  \defaultfontfeatures{Scale=MatchLowercase}
  \defaultfontfeatures[\rmfamily]{Ligatures=TeX,Scale=1}
\fi
\usepackage{lmodern}
\ifPDFTeX\else
  % xetex/luatex font selection
  \setmainfont[]{DejaVu Serif}
  \setsansfont[]{DejaVu Sans}
  \setmonofont[]{DejaVu Sans Mono}
\fi
% Use upquote if available, for straight quotes in verbatim environments
\IfFileExists{upquote.sty}{\usepackage{upquote}}{}
\IfFileExists{microtype.sty}{% use microtype if available
  \usepackage[]{microtype}
  \UseMicrotypeSet[protrusion]{basicmath} % disable protrusion for tt fonts
}{}
\makeatletter
\@ifundefined{KOMAClassName}{% if non-KOMA class
  \IfFileExists{parskip.sty}{%
    \usepackage{parskip}
  }{% else
    \setlength{\parindent}{0pt}
    \setlength{\parskip}{6pt plus 2pt minus 1pt}}
}{% if KOMA class
  \KOMAoptions{parskip=half}}
\makeatother
\usepackage{xcolor}
\usepackage[margin=0.5in]{geometry}
\usepackage{longtable,booktabs,array}
\usepackage{calc} % for calculating minipage widths
% Correct order of tables after \paragraph or \subparagraph
\usepackage{etoolbox}
\makeatletter
\patchcmd\longtable{\par}{\if@noskipsec\mbox{}\fi\par}{}{}
\makeatother
% Allow footnotes in longtable head/foot
\IfFileExists{footnotehyper.sty}{\usepackage{footnotehyper}}{\usepackage{footnote}}
\makesavenoteenv{longtable}
\setlength{\emergencystretch}{3em} % prevent overfull lines
\providecommand{\tightlist}{%
  \setlength{\itemsep}{0pt}\setlength{\parskip}{0pt}}
\setcounter{secnumdepth}{-\maxdimen} % remove section numbering
\ifLuaTeX
\usepackage[bidi=basic]{babel}
\else
\usepackage[bidi=default]{babel}
\fi
\babelprovide[main,import]{russian}
\ifPDFTeX
\else
\babelfont{rm}[]{DejaVu Serif}
\fi
% get rid of language-specific shorthands (see #6817):
\let\LanguageShortHands\languageshorthands
\def\languageshorthands#1{}
\ifLuaTeX
  \usepackage{selnolig}  % disable illegal ligatures
\fi
\usepackage{bookmark}
\IfFileExists{xurl.sty}{\usepackage{xurl}}{} % add URL line breaks if available
\urlstyle{same}
\hypersetup{
  pdflang={ru-RU},
  colorlinks=true,
  linkcolor={blue},
  filecolor={Maroon},
  citecolor={Blue},
  urlcolor={blue},
  pdfcreator={LaTeX via pandoc}}

\author{}
\date{}

\begin{document}

\section{1. Основные
параметры}\label{ux43eux441ux43dux43eux432ux43dux44bux435-ux43fux430ux440ux430ux43cux435ux442ux440ux44b}

Пусть:

\begin{itemize}
\tightlist
\item
  \(N\) --- количество игроков
\item
  \(M\) --- количество игровых месяцев
\item
  \(K_0\) --- начальный капитал игрока (по умолчанию 37 млн у.е.)
\item
  \(C\) --- полная стоимость строительства объекта
\item
  \(T\) --- срок строительства объекта в месяцах
\end{itemize}

Каждый шаг моделирования соответствует одному месяцу. Игра стартует в
марте (допущение, см. раздел 14).

\subsection{Параметры
объектов}\label{ux43fux430ux440ux430ux43cux435ux442ux440ux44b-ux43eux431ux44aux435ux43aux442ux43eux432}

Параметры объектов меняются в настройках.

\begin{longtable}[]{@{}
  >{\raggedright\arraybackslash}p{(\columnwidth - 8\tabcolsep) * \real{0.1579}}
  >{\raggedleft\arraybackslash}p{(\columnwidth - 8\tabcolsep) * \real{0.2105}}
  >{\raggedleft\arraybackslash}p{(\columnwidth - 8\tabcolsep) * \real{0.2105}}
  >{\raggedleft\arraybackslash}p{(\columnwidth - 8\tabcolsep) * \real{0.2105}}
  >{\raggedleft\arraybackslash}p{(\columnwidth - 8\tabcolsep) * \real{0.2105}}@{}}
\toprule\noalign{}
\begin{minipage}[b]{\linewidth}\raggedright
Объект
\end{minipage} & \begin{minipage}[b]{\linewidth}\raggedleft
Стоимость \(C\), млн
\end{minipage} & \begin{minipage}[b]{\linewidth}\raggedleft
Срок \(T\), мес.
\end{minipage} & \begin{minipage}[b]{\linewidth}\raggedleft
Квартир
\end{minipage} & \begin{minipage}[b]{\linewidth}\raggedleft
Площадь квартиры, кв. м
\end{minipage} \\
\midrule\noalign{}
\endhead
\bottomrule\noalign{}
\endlastfoot
Панельный дом & 8 & 7 & 60 & 50 \\
Монолитный дом & 10 & 9 & 60 & 55 \\
Кирпичный дом & 12 & 10 & 60 & 60 \\
Супермаркет & 2,5 & 5 & --- & --- \\
\end{longtable}

Себестоимость 1 кв. м жилья = стоимость дома / общая площадь дома.

\begin{center}\rule{0.5\linewidth}{0.5pt}\end{center}

\section{2. Затраты на
строительство}\label{ux437ux430ux442ux440ux430ux442ux44b-ux43dux430-ux441ux442ux440ux43eux438ux442ux435ux43bux44cux441ux442ux432ux43e}

Стоимость строительства расходуется равномерно в течение всего срока
строительства.

Ежемесячные затраты на один объект:

\[
C_{\text{month}} = \frac{C}{T}
\]

Пример: дом за 8 млн на 7 месяцев:

\[
\frac{8}{7} \approx 1{,}14
\]

млн у.е. в месяц.

Если строится несколько объектов, затраты за месяц складываются:

\[
C_{\text{build}}
=
\frac{C_1}{T_1}
+
\frac{C_2}{T_2}
+
\dots
\]

Если у игрока не хватает денег на платёж, стройка этого объекта в этом
месяце замораживается (затраты не списываются, срок сдвигается на
месяц).

\begin{center}\rule{0.5\linewidth}{0.5pt}\end{center}

\section{3. Спрос на
жильё}\label{ux441ux43fux440ux43eux441-ux43dux430-ux436ux438ux43bux44cux451}

\begin{quote}
\textbf{Примечание:} числа ниже --- допущения модели, см. раздел 14.
\end{quote}

Базовый спрос \(D_{\text{base}}(t)\) --- количество квартир в месяц для
каждого типа дома.

Он зависит от сезона:

\[
D_{\text{base}}(t)
=
25 \times (\text{коэффициент месяца})
\]

\begin{longtable}[]{@{}
  >{\raggedright\arraybackslash}p{(\columnwidth - 24\tabcolsep) * \real{0.0588}}
  >{\raggedleft\arraybackslash}p{(\columnwidth - 24\tabcolsep) * \real{0.0784}}
  >{\raggedleft\arraybackslash}p{(\columnwidth - 24\tabcolsep) * \real{0.0784}}
  >{\raggedleft\arraybackslash}p{(\columnwidth - 24\tabcolsep) * \real{0.0784}}
  >{\raggedleft\arraybackslash}p{(\columnwidth - 24\tabcolsep) * \real{0.0784}}
  >{\raggedleft\arraybackslash}p{(\columnwidth - 24\tabcolsep) * \real{0.0784}}
  >{\raggedleft\arraybackslash}p{(\columnwidth - 24\tabcolsep) * \real{0.0784}}
  >{\raggedleft\arraybackslash}p{(\columnwidth - 24\tabcolsep) * \real{0.0784}}
  >{\raggedleft\arraybackslash}p{(\columnwidth - 24\tabcolsep) * \real{0.0784}}
  >{\raggedleft\arraybackslash}p{(\columnwidth - 24\tabcolsep) * \real{0.0784}}
  >{\raggedleft\arraybackslash}p{(\columnwidth - 24\tabcolsep) * \real{0.0784}}
  >{\raggedleft\arraybackslash}p{(\columnwidth - 24\tabcolsep) * \real{0.0784}}
  >{\raggedleft\arraybackslash}p{(\columnwidth - 24\tabcolsep) * \real{0.0784}}@{}}
\toprule\noalign{}
\begin{minipage}[b]{\linewidth}\raggedright
Месяц
\end{minipage} & \begin{minipage}[b]{\linewidth}\raggedleft
Янв
\end{minipage} & \begin{minipage}[b]{\linewidth}\raggedleft
Фев
\end{minipage} & \begin{minipage}[b]{\linewidth}\raggedleft
Мар
\end{minipage} & \begin{minipage}[b]{\linewidth}\raggedleft
Апр
\end{minipage} & \begin{minipage}[b]{\linewidth}\raggedleft
Май
\end{minipage} & \begin{minipage}[b]{\linewidth}\raggedleft
Июн
\end{minipage} & \begin{minipage}[b]{\linewidth}\raggedleft
Июл
\end{minipage} & \begin{minipage}[b]{\linewidth}\raggedleft
Авг
\end{minipage} & \begin{minipage}[b]{\linewidth}\raggedleft
Сен
\end{minipage} & \begin{minipage}[b]{\linewidth}\raggedleft
Окт
\end{minipage} & \begin{minipage}[b]{\linewidth}\raggedleft
Ноя
\end{minipage} & \begin{minipage}[b]{\linewidth}\raggedleft
Дек
\end{minipage} \\
\midrule\noalign{}
\endhead
\bottomrule\noalign{}
\endlastfoot
Коэффициент спроса & 0,6 & 0,6 & 0,7 & 0,8 & 0,9 & 1,0 & 1,1 & 1,2 & 1,4
& 1,5 & 1,3 & 0,8 \\
\end{longtable}

Спрос растёт с весны и максимален осенью.

Супермаркеты в микрорайоне увеличивают спрос на жильё. Каждый
действующий супермаркет добавляет \(4\%\), каждый строящийся ---
\(2\%\):

\[
D(t)
=
D_{\text{base}}(t)
\times
\left(
1 + 0{,}04 \cdot S(t)
\right)
\]

где:

\begin{itemize}
\tightlist
\item
  \(D(t)\) --- общий спрос на жильё
\item
  \(S(t)\) --- число действующих супермаркетов плюс половина числа
  строящихся (у всех игроков вместе)
\end{itemize}

По формуле \(D(t)\) может получиться дробное число. Перед распределением
продаж текущая реализация округляет спрос вниз до целого числа квартир.

\begin{center}\rule{0.5\linewidth}{0.5pt}\end{center}

\section{4. Объём предложения
жилья}\label{ux43eux431ux44aux451ux43c-ux43fux440ux435ux434ux43bux43eux436ux435ux43dux438ux44f-ux436ux438ux43bux44cux44f}

Каждый игрок в начале месяца определяет объём жилья, который он
выставляет на продажу.

Выставить можно только доступное жильё: в готовом доме --- всё
непроданное, в строящемся --- часть, пропорциональную готовности дома.
Количество доступных квартир в строящемся доме определяется целочисленно:
\( \lfloor n \cdot t_{\text{built}}/T \rfloor \), где \(n\) --- число
квартир, \(t_{\text{built}}\) --- уже пройденные месяцы, \(T\) --- срок
строительства. Дробная часть квартиры отбрасывается.

\subsection{Единицы
измерения}\label{ux435ux434ux438ux43dux438ux446ux44b-ux438ux437ux43cux435ux440ux435ux43dux438ux44f}

Спрос и предложение считаются в \textbf{квартирах}, цена --- в у.е. за
\textbf{1 кв. м}.

Чтобы перейти к деньгам, нужна площадь квартиры \(a\) (зависит от типа
дома, см. таблицу в разделе 1).

Предложение игрока \(i\) (в квартирах):

\[
Q_i(t)
\]

Общее предложение:

\[
Q_{\text{total}}(t)
=
Q_1(t)
+
Q_2(t)
+
\dots
+
Q_N(t)
\]

\begin{center}\rule{0.5\linewidth}{0.5pt}\end{center}

\section{5. Влияние рекламы
жилья}\label{ux432ux43bux438ux44fux43dux438ux435-ux440ux435ux43aux43bux430ux43cux44b-ux436ux438ux43bux44cux44f}

Пусть \(A_i(t)\) --- затраты игрока \(i\) на рекламу жилья в тыс. у.е.

Каждая 1 тыс. рекламы увеличивает продажи на \(0{,}5\%\) в текущем и
следующем месяце:

\[
b_i
=
0{,}5\%
\times
\left(
\text{реклама в этом месяце}
+
\text{реклама в прошлом месяце}
\right)
\]

где \(b_i\) --- рекламный бонус игрока \(i\) к продажам жилья.

Пример: в прошлом месяце 2 тыс., в этом 4 тыс.:

\[
0{,}5\% \times 6 = 3\%
\]

\begin{center}\rule{0.5\linewidth}{0.5pt}\end{center}

\section{6. Распределение продаж
жилья}\label{ux440ux430ux441ux43fux440ux435ux434ux435ux43bux435ux43dux438ux435-ux43fux440ux43eux434ux430ux436-ux436ux438ux43bux44cux44f}

Продажа жилья идёт через единое риэлторское агентство.

Агентство работает отдельно для каждого типа дома:

\begin{enumerate}
\def\labelenumi{\arabic{enumi}.}
\item
  Если общее предложение не превышает спрос:

  \[
  Q_{\text{total}}(t) \leq D(t)
  \]

  всё выставленное жильё продаётся.
\item
  Если:

  \[
  Q_{\text{total}}(t) > D(t)
  \]

  предложения сортируются по цене от самой низкой, и спрос
  удовлетворяется по очереди: сначала самая дешёвая цена, пока хватает
  спроса.
\item
  Если цены одинаковые, остаток спроса делится между игроками
  пропорционально их предложению.
\item
  Предложения с ценой выше предельной не принимаются.

  Предельная цена:

  \[
  p_{\text{max}}
  =
  1{,}5
  \times
  p_{\text{fair}}
  \]

  где справедливая цена:

  \[
  p_{\text{fair}}
  =
  1{,}4
  \times
  c
  \]

  Здесь \(c\) --- себестоимость 1 кв. м.
\item
  Затем применяется реклама: базовое число проданных квартир
  умножается на \((1+b_i)\), но не больше выставленного объёма. Результат
  округляется вниз до целого числа квартир; поэтому небольшой рекламный
  бонус не всегда увеличивает фактическое количество продаж.
\end{enumerate}

\subsection{Пример}\label{ux43fux440ux438ux43cux435ux440}

Панельный дом, квартира 50 кв. м.

Спрос --- 100 квартир.

Игрок A предлагает 60 квартир по 2000 у.е. за кв. м.

Игрок B предлагает 70 квартир по 2200 у.е. за кв. м.

Игрок A продаёт:

\[
q_A = 60
\]

Остаётся спрос:

\[
100 - 60 = 40
\]

Игрок B продаёт:

\[
q_B = 40
\]

\begin{center}\rule{0.5\linewidth}{0.5pt}\end{center}

\section{7. Продажа жилья и
выручка}\label{ux43fux440ux43eux434ux430ux436ux430-ux436ux438ux43bux44cux44f-ux438-ux432ux44bux440ux443ux447ux43aux430}

Пусть:

\begin{itemize}
\tightlist
\item
  \(q_i(t)\) --- фактически проданное количество \textbf{квартир} игрока
  \(i\)
\item
  \(a\) --- площадь одной квартиры, кв. м (для данного типа дома)
\item
  \(p_i(t)\) --- цена \textbf{1 кв. м} жилья игрока \(i\)
\end{itemize}

Выручка от продажи жилья (если продаются дома разных типов, считается по
каждому типу и суммируется):

\[
R_{\text{house},i}(t)
=
q_i(t)
\cdot
a
\cdot
p_i(t)
\]

\subsection{Пример}\label{ux43fux440ux438ux43cux435ux440-1}

Продано 40 квартир по 50 кв. м по цене 2000 у.е. за кв. м:

\[
40 \cdot 50 \cdot 2000
=
4\,000\,000
\]

или:

\[
4\text{ млн у.е.}
\]

Выручка от продаж предыдущего месяца зачисляется игроку в начале
текущего месяца.

\begin{center}\rule{0.5\linewidth}{0.5pt}\end{center}

\section{8. Реклама
супермаркетов}\label{ux440ux435ux43aux43bux430ux43cux430-ux441ux443ux43fux435ux440ux43cux430ux440ux43aux435ux442ux43eux432}

Пусть \(B_i(t)\) --- затраты игрока \(i\) на рекламу товаров в
супермаркетах в тыс. у.е.

Реклама возможна, только если есть построенный супермаркет.

\(B_i\) --- общий рекламный бюджет игрока на все его супермаркеты:
надбавка \(w_i\) применяется к прибыли \textbf{каждого} его построенного
супермаркета, бюджет между ними не делится.

Каждые 0,5 тыс. рекламы увеличивают прибыль от продаж на \(3\%\).

\[
w_i
=
3\%
\times
\frac{B_i}{0{,}5}
\]

где \(w_i\) --- рекламная надбавка к прибыли супермаркетов игрока \(i\).

Пример: 2 тыс.:

\[
\frac{2}{0{,}5}
=
4
\]

блока по 0,5 тыс.

Следовательно:

\[
w_i
=
4 \times 3\%
=
12\%
\]

\begin{center}\rule{0.5\linewidth}{0.5pt}\end{center}

\section{9. Сезонность продаж в
супермаркетах}\label{ux441ux435ux437ux43eux43dux43dux43eux441ux442ux44c-ux43fux440ux43eux434ux430ux436-ux432-ux441ux443ux43fux435ux440ux43cux430ux440ux43aux435ux442ux430ux445}

Базовая прибыль \textbf{одного} действующего супермаркета в месяц ---
фиксированная сумма, не зависящая от его стоимости:

\[
P_{\text{base}}(t)
=
0{,}08\text{ млн}
\times
(\text{коэффициент месяца})
\]

\begin{longtable}[]{@{}
  >{\raggedright\arraybackslash}p{(\columnwidth - 24\tabcolsep) * \real{0.0588}}
  >{\raggedleft\arraybackslash}p{(\columnwidth - 24\tabcolsep) * \real{0.0784}}
  >{\raggedleft\arraybackslash}p{(\columnwidth - 24\tabcolsep) * \real{0.0784}}
  >{\raggedleft\arraybackslash}p{(\columnwidth - 24\tabcolsep) * \real{0.0784}}
  >{\raggedleft\arraybackslash}p{(\columnwidth - 24\tabcolsep) * \real{0.0784}}
  >{\raggedleft\arraybackslash}p{(\columnwidth - 24\tabcolsep) * \real{0.0784}}
  >{\raggedleft\arraybackslash}p{(\columnwidth - 24\tabcolsep) * \real{0.0784}}
  >{\raggedleft\arraybackslash}p{(\columnwidth - 24\tabcolsep) * \real{0.0784}}
  >{\raggedleft\arraybackslash}p{(\columnwidth - 24\tabcolsep) * \real{0.0784}}
  >{\raggedleft\arraybackslash}p{(\columnwidth - 24\tabcolsep) * \real{0.0784}}
  >{\raggedleft\arraybackslash}p{(\columnwidth - 24\tabcolsep) * \real{0.0784}}
  >{\raggedleft\arraybackslash}p{(\columnwidth - 24\tabcolsep) * \real{0.0784}}
  >{\raggedleft\arraybackslash}p{(\columnwidth - 24\tabcolsep) * \real{0.0784}}@{}}
\toprule\noalign{}
\begin{minipage}[b]{\linewidth}\raggedright
Месяц
\end{minipage} & \begin{minipage}[b]{\linewidth}\raggedleft
Янв
\end{minipage} & \begin{minipage}[b]{\linewidth}\raggedleft
Фев
\end{minipage} & \begin{minipage}[b]{\linewidth}\raggedleft
Мар
\end{minipage} & \begin{minipage}[b]{\linewidth}\raggedleft
Апр
\end{minipage} & \begin{minipage}[b]{\linewidth}\raggedleft
Май
\end{minipage} & \begin{minipage}[b]{\linewidth}\raggedleft
Июн
\end{minipage} & \begin{minipage}[b]{\linewidth}\raggedleft
Июл
\end{minipage} & \begin{minipage}[b]{\linewidth}\raggedleft
Авг
\end{minipage} & \begin{minipage}[b]{\linewidth}\raggedleft
Сен
\end{minipage} & \begin{minipage}[b]{\linewidth}\raggedleft
Окт
\end{minipage} & \begin{minipage}[b]{\linewidth}\raggedleft
Ноя
\end{minipage} & \begin{minipage}[b]{\linewidth}\raggedleft
Дек
\end{minipage} \\
\midrule\noalign{}
\endhead
\bottomrule\noalign{}
\endlastfoot
Коэффициент продаж & 1,5 & 1,3 & 1,1 & 0,9 & 0,8 & 0,8 & 0,8 & 0,9 & 1,0
& 1,1 & 1,2 & 1,4 \\
\end{longtable}

Продажи растут с осени и максимальны зимой.

Дома в микрорайоне увеличивают продажи супермаркетов: каждый дом даёт
\(+2\%\) (строящийся --- \(+1\%\)).

Итоговая прибыль одного супермаркета:

\[
P
=
P_{\text{base}}(t)
\times
\left(
1 + 0{,}02 \cdot H
\right)
\times
\left(
1 + w_i
\right)
\]

где:

\begin{itemize}
\tightlist
\item
  \(H\) --- число готовых домов плюс половина строящихся (у всех игроков
  вместе)
\item
  \(w_i\) --- надбавка за рекламу из раздела 8 (например, \(0{,}12\) для
  \(+12\%\))
\end{itemize}

Прибыль за прошлый месяц зачисляется в начале текущего.

\begin{center}\rule{0.5\linewidth}{0.5pt}\end{center}

\section{10. Стоимость объектов при расчёте
капитала}\label{ux441ux442ux43eux438ux43cux43eux441ux442ux44c-ux43eux431ux44aux435ux43aux442ux43eux432-ux43fux440ux438-ux440ux430ux441ux447ux451ux442ux435-ux43aux430ux43fux438ux442ux430ux43bux430}

Каждый объект учитывается \textbf{только в одной} составляющей капитала:

\begin{itemize}
\tightlist
\item
  \(C_{\text{construction}}\) --- только \textbf{незавершённые} объекты.
  Для строящихся супермаркетов учитывается вся потраченная сумма; для
  строящегося дома в текущей реализации из потраченной суммы вычитается
  себестоимость уже проданных квартир;
\item
  \(C_{\text{housing}}\) --- только \textbf{непроданное жильё в
  построенных домах}: площадь непроданных квартир × себестоимость 1 кв.
  м;
\item
  \(C_{\text{markets}}\) --- только \textbf{построенные супермаркеты}:
\end{itemize}

\[
C_{\text{markets}}
=
1{,}6
\cdot
C_{\text{spent}}
\]

где \(C_{\text{spent}}\) --- сумма, потраченная на строительство
супермаркета.

\begin{center}\rule{0.5\linewidth}{0.5pt}\end{center}

\section{11. Общий капитал
игрока}\label{ux43eux431ux449ux438ux439-ux43aux430ux43fux438ux442ux430ux43b-ux438ux433ux440ux43eux43aux430}

Обозначим:

\begin{itemize}
\tightlist
\item
  \(C_{\text{construction}}\) --- стоимость строящихся (незавершённых)
  объектов
\item
  \(C_{\text{housing}}\) --- стоимость непроданного жилья в построенных
  домах
\item
  \(C_{\text{markets}}\) --- стоимость построенных супермаркетов
\item
  \(Cash\) --- денежные средства игрока
\end{itemize}

Тогда общий капитал игрока:

\[
K
=
C_{\text{construction}}
+
C_{\text{housing}}
+
C_{\text{markets}}
+
Cash
\]

Денежные средства: начальный капитал минус все расходы (строительство,
реклама) плюс все доходы (продажи жилья, прибыль супермаркетов).

Баланс не может быть отрицательным.

Проданное жильё учитывается только в \(Cash\) (как полученные деньги) и
в стоимость объектов больше не входит.

\begin{quote}
\textbf{Допущение реализации, отличающееся от исходного Word-задания:}
квартиры, проданные в ещё строящемся доме, уже превращены в деньги, поэтому
из «потраченной суммы» этого дома вычитается их себестоимость --- иначе одна
и та же часть дома была бы учтена дважды (в деньгах и в стоимости
строящегося объекта). В исходном задании стоимость незавершённого объекта
описана как вся сумма, уже потраченная на его строительство.
\end{quote}

\begin{center}\rule{0.5\linewidth}{0.5pt}\end{center}

\section{12. Условие
победы}\label{ux443ux441ux43bux43eux432ux438ux435-ux43fux43eux431ux435ux434ux44b}

После выполнения \(M\) игровых шагов для каждого игрока рассчитывается
итоговый капитал:

\[
K_i(M)
\]

\begin{quote}
\textbf{Порядок начисления в текущей реализации:} выручка от продаж
жилья месяца \(M\) зачисляется в \(Cash\) при подведении итогов; настройки
для отключения этого поведения в \texttt{Config} нет. Прибыль супермаркетов,
рассчитанная по результатам месяца \(M\), остаётся отложенной и не входит
в итоговый капитал, поскольку в обычном цикле она зачисляется в начале
следующего месяца.
\end{quote}

Победителем становится игрок с максимальным капиталом:

\[
K_{\text{winner}}
=
\max_i K_i(M)
\]

При равенстве капиталов объявляется ничья.

\begin{center}\rule{0.5\linewidth}{0.5pt}\end{center}

\section{13. Последовательность расчётов за
месяц}\label{ux43fux43eux441ux43bux435ux434ux43eux432ux430ux442ux435ux43bux44cux43dux43eux441ux442ux44c-ux440ux430ux441ux447ux451ux442ux43eux432-ux437ux430-ux43cux435ux441ux44fux446}

\subsection{Начальная
фаза}\label{ux43dux430ux447ux430ux43bux44cux43dux430ux44f-ux444ux430ux437ux430}

Перед месяцем 1 каждый игрок получает стартовый капитал и решает, в
какие объекты вложить деньги: выбирает, какие стройки начать.

Первые платежи за эти стройки списываются в шаге 1 месяца 1.

\subsection{Новые стройки в ходе
игры}\label{ux43dux43eux432ux44bux435-ux441ux442ux440ux43eux439ux43aux438-ux432-ux445ux43eux434ux435-ux438ux433ux440ux44b}

Решение начать новую стройку принимается в шаге 4 месяца \(t\).

Первый платёж списывается в шаге 1 месяца \(t+1\).

\subsection{Расчёты за
месяц}\label{ux440ux430ux441ux447ux451ux442ux44b-ux437ux430-ux43cux435ux441ux44fux446}

Для каждого месяца \(t\) выполняется следующая последовательность:

\begin{enumerate}
\def\labelenumi{\arabic{enumi}.}
\tightlist
\item
  рассчитываются и списываются затраты на строительство;
\item
  начисляется выручка от продаж жилья предыдущего месяца;
\item
  начисляется прибыль супермаркетов;
\item
  игроки принимают решения на месяц: новые стройки, предложение жилья;
\item
  устанавливаются цены;
\item
  учитываются затраты на рекламу жилья;
\item
  учитываются затраты на рекламу в супермаркетах;
\item
  риэлторское агентство распределяет спрос и определяет объём продаж;
\item
  обновляется состояние объектов и капитал игроков.
\end{enumerate}

После последнего месяца рассчитывается итоговый капитал всех игроков и
определяется победитель.

\begin{center}\rule{0.5\linewidth}{0.5pt}\end{center}

\section{14. Допущения
модели}\label{ux434ux43eux43fux443ux449ux435ux43dux438ux44f-ux43cux43eux434ux435ux43bux438}

Значения ниже не заданы в ТЗ и приняты как допущения модели. Их можно
менять в настройках.

\begin{longtable}[]{@{}
  >{\raggedright\arraybackslash}p{(\columnwidth - 2\tabcolsep) * \real{0.5000}}
  >{\raggedright\arraybackslash}p{(\columnwidth - 2\tabcolsep) * \real{0.5000}}@{}}
\toprule\noalign{}
\begin{minipage}[b]{\linewidth}\raggedright
Допущение
\end{minipage} & \begin{minipage}[b]{\linewidth}\raggedright
Значение по умолчанию
\end{minipage} \\
\midrule\noalign{}
\endhead
\bottomrule\noalign{}
\endlastfoot
Месяц старта игры & март \\
Сезонные коэффициенты спроса и продаж & таблицы в разделах 3 и 9 \\
Базовый спрос & 25 квартир в месяц на каждый тип дома \\
Влияние супермаркетов на спрос & +4 \% за действующий, +2 \% за
строящийся \\
Влияние домов на продажи супермаркетов & +2 \% за готовый, +1 \% за
строящийся \\
Базовая прибыль супермаркета & 0,08 млн в месяц \\
Параметры монолитного и кирпичного домов, число квартир и площади &
таблица раздела 1 \\
Предельная цена & 1,5 × 1,4 × себестоимость кв. м \\
Нехватка денег на платёж & стройка замораживается на месяц \\
Реклама жилья & привлекает покупателей сверх спроса, но не больше
выставленного объёма \\
Выручка от продаж жилья в месяце \(M\) & зачисляется при подведении итогов \\
Прибыль супермаркетов за месяц \(M\) & остаётся отложенной и не входит в итоговый капитал \\
Жильё, проданное в строящемся доме & вычитается из стоимости этого дома
(раздел 11) \\
Ничья & объявляется ничья \\
\end{longtable}

\end{document}
